\chapter{Model Results}
% OUTLINE ONLY -- replace the comments with text.
This chapter will present what type of model best predicts the data, and what amount of states
best describe the data. We will then do a analysis of the behaviour of our choosen model
%Maybe just say that we will chose AR(1)-HMM(B,1) for 3 or 4 states

\section{Model type}
As desribed in experimental setup, we fit 4 models with a increasing order of complexity. \
Each increase adds a new layer to the structure of how the model models the dynamics in the room.\

\begin{table}[ht]
\centering
\begin{adjustbox}{max width=\textwidth}
\begin{tabular}{lcccc}
\hline
States & HMM(1) & HMM(b,1) & AR(1)-HMM(1) & AR(1)-HMM(b,1) \\
\hline
3 & 333.62 / 19.97 & 325.57 / 19.27 & 151.86 / 5.74 & 150.64 / 5.65 \\
4 & 285.94 / 14.69 & 281.95 / 13.66 & 151.94 / 5.57 & 151.23 / 5.94 \\
5 & 244.04 / 12.96 & 243.44 / 11.35 & 151.00 / 5.50 & 151.94 / 5.90 \\
6 & 249.13 / 12.98 & 203.23 / 9.95 & 150.98 / 5.59 & 151.69 / 5.67 \\
7 & 215.36 / 10.45 & 204.17 / 9.26 & 151.80 / 5.64 & 151.05 / 5.57 \\
\hline
\end{tabular}
\end{adjustbox}
\caption{Validation RMSE (ppm) / MAPE (\%) of the 1-step-ahead forecasts for each model and number of states.}
\label{tab:predictive_metrics_room_015_calibrated_anchor_random_search_val_h1}
\end{table}

\begin{table}[ht]
\centering
\begin{adjustbox}{max width=\textwidth}
\begin{tabular}{lcccc}
\hline
States & HMM(1) & HMM(b,1) & AR(1)-HMM(1) & AR(1)-HMM(b,1) \\
\hline
3 & 507.40 / 39.66 & 405.21 / 28.74 & 414.21 / 27.13 & 349.72 / 19.95 \\
4 & 495.06 / 34.38 & 389.70 / 25.67 & 419.22 / 24.73 & 346.95 / 19.78 \\
5 & 483.60 / 31.96 & 365.20 / 22.88 & 408.45 / 23.53 & 351.82 / 23.60 \\
6 & 481.57 / 32.30 & 356.14 / 21.55 & 421.94 / 30.88 & 345.79 / 20.67 \\
7 & 478.87 / 29.90 & 357.94 / 20.65 & 426.11 / 30.97 & 360.14 / 18.75 \\
\hline
\end{tabular}
\end{adjustbox}
\caption{Validation RMSE (ppm) / MAPE (\%) of the 8-steps-ahead forecasts for each model and number of states.}
\label{tab:predictive_metrics_room_015_calibrated_anchor_random_search_val_h8}
\end{table}





\section{Choosing the Number of States}
% Table of LL / AIC / BIC per K for each model, next to validation RMSE / R^2.
% Focus: argue for one K. Information criteria tend to keep rising with K, so
% use predictive performance + interpretability (do the extra states mean anything?).

\section{Model Comparison}

\subsection{In-sample Fit}
% LL, AIC, BIC across the model ladder for the chosen K.
% Focus: what each added component (AR(1) emission, covariates in transition,
% shared vs state-specific beta) buys, and whether it is worth its parameters.

\subsection{Out-of-sample Forecasting}
% RMSE, R^2, MAPE per forecast horizon on validation and test (test_metrics.csv).
% Focus: this is the main ranking criterion. Show metric vs horizon plot,
% compare against a simple baseline (e.g. persistence), note where models diverge.

\section{The Selected Model}
\subsection{Estimated Parameters}
% Table of state means, sigmas, phi (AR coefficient) and the transition matrix / betas.
% Focus: physical interpretation -- link phi to ventilation/decay rate (OU, Ch. 3),
% mu to occupancy levels; comment on parameter uncertainty if available.

\subsection{Decoded States}
% Decoded/smoothed state sequence over a representative week (decode_states).
% Focus: do states line up with occupied / empty / ventilation periods and with
% the working-hours / off-day covariates? Use transition_plots for the covariate effect.

\subsection{Forecast Behaviour}
% First-week forecast plots and the state confusion matrix (predicted vs smoothed state).
% Focus: how well the model anticipates regime changes, not just level tracking.

\section{Model Diagnostics}
% Pseudo-residuals: QQ-plot, histogram, ACF, residuals over time.
% Focus: are residuals approx. N(0,1) and uncorrelated? Name remaining misfit
% (e.g. heavy tails, unmodelled seasonality) -- this feeds the Discussion.

\section{Robustness}
% Optional: results on a second room (room_009) and/or the other data tags.
% Focus: does the model ranking and chosen K hold up, or is it room-specific?

\section{Summary}
% A few bullet-point conclusions: chosen model and K, its predictive gain,
% key physical interpretation, main shortcomings carried into Ch. 6/7.
